El propósito de este trabajo es estudiar el concepto de sucesión equidistribuida módulo 1 y probar algunos criterios que garantizan que una sucesión tiene esta propiedad, entre ellos el teorema de Weyl. Una sucesión $(x_n)_{n\geq1}$ módulo 1 consiste en observar solamente la parte fraccionaria de sus términos $\langle x_n\rangle$. Analizaremos cómo se distribuyen estos valores en el intervalo $[0,1)$, y en particular, cuando se distribuyen uniformemente en dicho intervalo. A lo largo de la memoria, desarrollaremos herramientas eficaces para demostrar esta propiedad.

El trabajo se divide en cuatro capítulos. En el primero de ellos, introduciremos la noción de los números reales módulo 1 y el espacio cociente $\mathbb{T}=\mathbb{R}/\mathbb{Z}$, conocido como el toro unidimensional, y que se identifica con el intervalo $[0,1)$ con sus extremos conectados. A continuación, definiremos el concepto de sucesión equidistribuida modulo 1 y demostraremos algunas propiedades fundamentales, acompañadas de ejemplos ilustrativos. Cerraremos este capítulo demostrando el teorema de equidistribución de Weyl, que establece la equidistribución de la sucesión $(n\alpha)$ 
para todo $\alpha$ irracional. Para ello, haremos uso de la aproximación de funciones continuas mediante polinomios trigonométricos.

El segundo capítulo se centra en desarrollar las principales herramientas para verificar la equidistribución de una sucesión. En primer lugar, presentamos la caracterización funcional, según la cual una sucesión es equidistribuida módulo 1 si y solo si, para cualquier función Riemann integrable la media de la función evaluada sobre los valores de la sucesión converge a la integral de dicha función en $\mathbb{T}$.
La segunda  es el criterio de Weyl. Este es uno de los métodos más importante, por no decir el más importante, para comprobar la equidistribución. Reduce el problema al cálculo asintótico de sumas de exponenciales complejas. Como aplicaciones se demuestra nuevamente la equidistribución de $(n\alpha)$ y se estudian ejemplos adicionales incluyendo sucesiones que no son uniformemente distribuidas $(\sin n)$.

%\footnote{{\color{blue}   Aqu\'i tienes la opci\'on de escribir alg\'un enunciado o alguna f\'ormula, de modo que el lector (tribunal) tenga m\'as claro a qu\'e te refieres...}}

En el tercer capítulo, se exploran técnicas de acotación para las sumas de exponenciales. Utilizaremos la fórmula de sumación de Euler. Esta nos permite reescribir las sumas discretas en términos de integrales las cuales pueden ser estimadas con precisión. De este modo, al acotarlas, podremos verificar la equidistribución mediante el criterio de Weyl. Demostraremos que la sucesión $(a\log n)$ no es equidistribuida módulo 1, mientras que sucesiones de la forma $(an^\sigma)$ con $0<\sigma<1$, sí lo son. El capítulo termina con la demostración del teorema de Fejér, que garantiza la equidistribución de una sucesión $(g(n))$ siempre que la función $g$ y su derivada $g'$ cumplan ciertas condiciones de crecimiento, límite y monotonía.

Para finalizar, estudiamos el método de las diferencias de Van der Corput. Este método se basa en la desigualdad de Van der Corput, que relaciona las sumas de exponenciales con las correspondientes sumas de sus diferencias. A partir de esta desigualdad, demostramos que una sucesión es equidistribuida si también lo son sus sucesiones de diferencias. Como aplicación principal, probamos que toda sucesión definida por un polinomio con al menos un coeficiente irracional no constante es equidistribuida modulo 1. Además, el método de las diferencias de Van der Corput permite extender el teorema de Fejér a funciones más generales, siempre que sus derivadas de orden superior cumplan las condiciones de crecimiento adecuadas.

